% ─────────────────────────────────────────────────────────────────────────────
\section{Final System Architecture and Design}
\label{sec:architecture}
% ─────────────────────────────────────────────────────────────────────────────

The final system architecture is a two-stage hybrid pipeline designed to combine the speed of
a specialized thermal detector with the reasoning depth of an open vision-language model.
This design was chosen because a single general-purpose model could not simultaneously
deliver reliable thermal target localization, low-latency scanning, and semantically useful
interpretation.

\subsection{Architecture Overview}

\begin{figure}[H]
  \centering
  \includegraphics[width=0.98\textwidth]{system_pipeline.png}
  \caption{High-level system pipeline used in the final prototype.}
\end{figure}

The architecture begins with a thermal frame, applies full-frame detection, isolates each
candidate region, performs target-focused reasoning, and then fuses the results into
structured output. This decomposition reflects the key engineering principle of the project:
\emph{detect first, interpret second}.

\subsection{Two-Stage ``Gözcü + Analist'' Pipeline}

The two main stages are:

\begin{table}[H]
  \centering
  \caption{Primary roles of the two-stage pipeline}
  \begin{tabularx}{\textwidth}{@{}p{3cm}p{3.2cm}X@{}}
    \toprule
    Stage & Component & Responsibility \\
    \midrule
    Stage 1 & Gözcü (Scout) & Scan the full scene quickly, detect objects of interest,
    and return bounding boxes, class labels, and confidence values. \\
    Stage 2 & Analist (Analyst) & Receive cropped target regions and infer deeper semantic
    meaning such as likely object type, possible engine activity, carried-object clues, or
    qualitative threat relevance. \\
    \bottomrule
  \end{tabularx}
\end{table}

This role separation is important. The detector answers ``where is the object?'' and ``what
is the coarse class?'' The reasoner answers ``what does this object likely mean in context?''

\subsection{Subsystem Decomposition}

The final architecture can be understood as five subsystems.

\begin{itemize}[nosep]
  \item \textbf{Data acquisition:} Accept local thermal or RGB streams, static images, or
    frame sequences without requiring external network services.
  \item \textbf{Detection:} Run a thermal-adapted YOLOv8 model on the full frame and
    generate detections suitable for downstream processing.
  \item \textbf{ROI processing:} Crop and optionally enhance the detected object region so
    the second stage focuses on the target instead of the whole scene.
  \item \textbf{Semantic reasoning:} Use Qwen2-VL with a prompt-driven analysis workflow
    to produce richer descriptions and infer status-oriented details.
  \item \textbf{Fusion and presentation:} Merge detector and reasoner outputs into JSON,
    logs, and local interface elements for operator review.
\end{itemize}

\subsection{Low-Level Deployment Packages}

Beyond the logical architecture, the implementation protocol was organized into deployment-
oriented packages so that hardware-specific concerns could be isolated from application-level
reasoning logic.

\begin{itemize}[nosep]
  \item \textbf{Containerization package:} Encapsulates the runtime environment, model
    dependencies, and platform-specific libraries inside a reproducible Docker image for
    Jetson-class deployment.
  \item \textbf{Optimization package:} Applies runtime acceleration techniques such as
    TensorRT or Torch-TensorRT compilation, layer fusion where applicable, and reduced-
    precision execution paths for edge hardware.
  \item \textbf{Inference pipeline package:} Coordinates frame ingestion, detector
    execution, region cropping, reasoning invocation, and result fusion under one bounded
    runtime flow.
  \item \textbf{Logging and presentation package:} Stores structured JSON outputs locally
    and exposes operator-facing summaries, alerts, or dashboard elements without external
    services.
\end{itemize}

This decomposition is valuable because it separates \emph{what} the system does from
\emph{how} it is deployed. The detector and reasoner remain conceptual building blocks, but
the deployment packages define how those blocks are actually started, optimized, and
operated on edge devices.

\subsection{Secure Deployment in a Closed Network}

\begin{figure}[H]
  \centering
  \includegraphics[width=0.98\textwidth]{secure_deployment.png}
  \caption{Closed-network deployment concept with local processing and no external
  connectivity.}
\end{figure}

The deployment concept keeps image ingestion, model inference, result fusion, and storage
within the same local environment. This provides three advantages:

\begin{itemize}[nosep]
  \item sensitive data does not leave the operational network,
  \item latency is reduced because inference is local,
  \item the system remains usable in environments where internet access is unavailable or
    intentionally disabled.
\end{itemize}

At low level, this also affects device integration. Video ingress must be available from a local
camera or local network stream, model binaries must reside on the target device or an
authorized internal store, and local visualization inside a containerized runtime may require
explicit device and display mapping between host and container.

\subsection{Control Flow}

The internal software flow of one processed frame is summarized below.

\begin{enumerate}[nosep]
  \item Acquire a frame from a local thermal source.
  \item Run the detector and obtain a list of candidate objects.
  \item For each detection, crop the relevant region of interest.
  \item Send the cropped region and an analysis prompt to the VLM.
  \item Receive qualitative reasoning output.
  \item Merge detector data and reasoner data into a structured record.
  \item Display and/or store the final result locally.
\end{enumerate}

\subsection{Interface Contracts}

The system is easier to maintain when each stage exposes a narrow and typed interface. The
final implementation protocol therefore treated component boundaries as explicit contracts,
even when the internal implementations evolved.

\begin{itemize}[nosep]
  \item \textbf{\texttt{ingest\_frame(source)}}: Returns a raw frame buffer from a local
    stream, file, or camera source.
  \item \textbf{\texttt{detect\_objects(frame)}}: Accepts a preprocessed frame and returns
    bounding boxes, class labels, and confidence values.
  \item \textbf{\texttt{crop\_regions(frame, detections)}}: Converts detector outputs into
    isolated regions of interest suitable for second-stage analysis.
  \item \textbf{\texttt{analyze\_region(region, prompt)}}: Produces a structured semantic
    description for one target crop.
  \item \textbf{\texttt{fuse\_results(detections, analyses)}}: Merges low-level detector
    data and higher-level reasoning output into a final JSON record.
\end{itemize}

The exact tensor dimensions are implementation-dependent, but the important engineering
rule is that each boundary must preserve shape, data type, and memory ownership clearly
enough to avoid runtime mismatch between preprocessing, inference, and fusion layers.

\subsection{Structured Output Format}

The final output is intentionally machine-readable so that it can be consumed by dashboards,
logs, or future command-and-control integrations.

\begin{lstlisting}[language={},caption={Illustrative structured output produced by the fusion layer}]
{
  "objects": [
    {
      "type": "armed_person",
      "bounding_box": [x1, y1, x2, y2],
      "confidence": 0.92,
      "analysis": {
        "weapon_type": "rifle",
        "risk_level": "high",
        "explanation": "Thermal silhouette suggests a long object consistent with a rifle."
      }
    },
    {
      "type": "vehicle",
      "bounding_box": [x3, y3, x4, y4],
      "confidence": 0.85,
      "analysis": {
        "vehicle_category": "pickup_truck",
        "engine_status": "hot",
        "explanation": "Front thermal concentration suggests recent engine activity."
      }
    }
  ]
}
\end{lstlisting}

\subsection{Design Rationale}

The final design was selected because it balances feasibility and ambition.

\begin{itemize}[nosep]
  \item A standalone VLM is too weak at thermal scene scanning.
  \item A standalone detector is too shallow for contextual interpretation.
  \item A two-stage pipeline is modular, explainable, and easier to improve incrementally.
  \item Local JSON output supports both immediate use and later auditability.
\end{itemize}

The result is not merely a collection of models, but a deliberately engineered workflow that
turns thermal pixels into structured, reviewable intelligence.
